\scanpage{113}
\deleted{et comme $s$ non lié à $w$ (sinon $\{s,u,v\}$ et $\{s,u,w\}$ \unsure{seraient|seront} des triangles), ni $t$ à $v$, on aura $r\neq v,w$. Donc $\sigma_{\underline\sigma}$ et $\sigma_{\underline\tau}$ sont l'identité sur $r$, donc \deleted{aussi sur $r$} $U=\sigma_{\underline\sigma}\sigma_{\underline\tau}$ aussi.}

\deleted{Donc, si \deleted{\unsure{Soit}} $\underline t,\underline t'\in\widetilde T$ sont deux triangles supérieurs qui coïncident en \unsure{le|au} sommet $c$ (\unsure{soit|sur} $r$) et qui ont \deleted{tels que} les sommets $s,t$ au dessus de $a$, -- alors, si on peut trouver $u\in\widetilde E_b$ lié \unsure{à la} fois à $s$ et $t$, alors}

\figure{\begin{tikzpicture}[x=0.016cm,y=-0.016cm,line cap=round,line join=round]
\draw (201,760)--(278,629);
\draw (285,627)--(340,757);
\draw (201,760)--(330,770);
\draw (305,411)--(443,500);
\draw (304,475)--(392,622);
\fill (305,411) circle[radius=4];
\fill (443,500) circle[radius=4];
\fill (304,475) circle[radius=4];
\fill (392,622) circle[radius=4];
\fill (247,469) circle[radius=2];
\fill (240,518) circle[radius=2];
\node[inner sep=1pt,font=\scriptsize] at (269,599) {$a$};
\node[inner sep=1pt,font=\scriptsize] at (377,769) {$b$};
\node[inner sep=1pt,font=\scriptsize] at (164,778) {$c$};
\node[inner sep=1pt,font=\scriptsize] at (270,455) {$t$};
\node[inner sep=1pt,font=\scriptsize] at (262,524) {$s$};
\node[inner sep=1pt,font=\scriptsize] at (476,710) {$u$};
\end{tikzpicture}}

% archive sheet number: 112
